\documentclass[Markos_Delaportas_1316.tex]{subfiles}

\begin{document}    
    
    \section{System Model and Practical Challenges in \texorpdfstring{\gls{ota}}{ }}
        A system that takes advantage of \gls{ota} can be modeled in different
        ways, also depending on the use case and environment. Most layouts involve
        a number of \gls{ed}s that transmit to a \gls{fc} or \gls{es}. The channel 
        model is also an important aspect, as it can be affected by fading, noise, 
        and path loss. In addition to the system model, there are also practical 
        challenges that can affect the reliability of \gls{ota}, such as synchronization 
        errors and waveform design. To enhance the performance of
        \gls{ota}, optimization strategies can be employed, such as power management
        and transceiver design.

    \subsection{System Architecture and Channel Model}
        In this section the mathematical model of a network consisting of three
        \gls{ed}s will be presented. These devices are deployed at an agricultural
        field and take advantage of \gls{ota} in order to transmit to the \gls{fc},
        by using \gls{pam}. The physical environment has multipath characteristics
        and an impulse response with a span of three taps.

    \begin{itemize}
        \item \textit{Modeling the Network}: Description of \gls{ed}s
            transmitting to an \gls{fc}, possibly involving multiple antennas
            (\gls{mimo})\cite{sahin_survey_2023,yang_federated_2020,
            hellstrom_optimal_2023,cao_optimized_2022,zhu_over--air_2024}.

            \begin{itemize}
                \item \textbf{Sensor 1}: $x_1[n] = [2, -1, 1]$
                \item \textbf{Sensor 2}: $x_1[n] = [1, 2, -1]$
            \end{itemize}

        \item \textit{Wireless Channel Model}: Consideration of
            channel impairments, including
            fading (e.g., Rayleigh fading), noise (\gls{awgn}), and path
            loss\cite{zhu_one-bit_2021,evgenidis_waveform_2025,cao_optimized_2022}.

            The channel impulse response is given as: $h[n] = [1.0, 0.5, -0.2]$

            (\textit{Here zero noise is assumed})
    \end{itemize}

    To calculate the ideal target function of the \gls{ota} aggregated sequence, $s[n]$,
    the ideal flat channel without multipath is assumed as: $h_{ideal}[n] = [1]$.
    \begin{equation}
        \label{eq:IdealTargetFunc}
        \begin{aligned}
            s[n] & = x_1[n] + x_2[n] \Rightarrow \\
            s[n] & = [3, 1, 0]
        \end{aligned}
    \end{equation}

    \subsection{Impact of Practical Impairments on OAC Reliability}
        In order to calculate the received signal with \gls{isi} at the \gls{fc},
        the linear property of the channel can be taken advantage of.

        \begin{equation}
            \label{eq:ActualTargetFunc}
            \begin{aligned}
                y[n] & = s[n] \ast h[n] \Rightarrow \\
                y[n] & = \sum_{m=0}^{n} s[m] h[n - m] \Rightarrow \\
                y[n] & = [3.0, 2.5, -0.1, -0.2, 0];
            \end{aligned}
        \end{equation}

        It is obvious that because of the \gls{isi} the three symbols of 
        \ref{eq:IdealTargetFunc} are spread to the five symbols of \ref{eq:ActualTargetFunc}.    

        To calculate the \gls{mse} at each sample the comparison of the actual received
        signal to the ideal target function shall be made.

        \begin{equation}
            \label{eq:mse}
            \begin{aligned}
                e[n] & = y[n] - n[n] \Rightarrow \\
                e[n] & = [0.0, 1.5, -0.1, -0.2, 0.0] \\
                MSE & = \frac{e[n]}{5} \Rightarrow \\
                MSE & = 0.46
            \end{aligned}
        \end{equation}

    \subsubsection{Synchronization Errors}
    \textit{Time, Frequency, and Phase Offsets (TO, CFO, PO)}: The detrimental
    impact of imperfect time-of-arrivals and phase alignment on signal
    superposition and computation
    accuracy\cite{sahin_survey_2023,evgenidis_waveform_2025,hellstrom_optimal_2023}.

    \begin{lstlisting}
        % This is a comment
        x = linspace(0, 2*pi, 100);
        y = sin(x) + 0.5 * randn(size(x)); % Numbers like 0.5 will be in Math Font
        plot(x, y);
    \end{lstlisting}

    \subsubsection{Channel State Information (CSI) Imperfection}
    The need for accurate and fresh \gls{csi} at transmitters or
    receivers for many precoding/equalization techniques
    \cite{sahin_survey_2023,yang_federated_2020,evgenidis_waveform_2025}.
    \subsubsection{Waveform and Filter Design}
    Addressing time sampling error and intersymbol interference \gls{isi}
    inherent in digital communication transceivers when applied to \gls{ota}
    \cite{evgenidis_waveform_2025}.
    \subsection{Optimization Strategies to Enhance OAC Performance}
    \begin{itemize}
        \item \textit{Power Management}: Strategies to mitigate
            aggregation errors, such as
            optimizing power control policies to minimize Mean Squared Error
            \gls{mse} or maximize convergence speed
            \cite{sahin_survey_2023,cao_optimized_2022,evgenidis_waveform_2025}.
        \item \textit{Transceiver Design}: Joint design of
            pre-processing functions,
            precoders/beamformers, and post-processing functions (decoders) to
            match the channel to the desired function and minimize
            distortion \cite{goldenbaum_robust_2013,hellstrom_optimal_2023,
            sahin_survey_2023,yang_federated_2020}.
    \end{itemize}

    \ifSubfilesClassLoaded{
        \clearpage
        \printbibliography
    }{}

\end{document}